\section{Procedural and Workflow Understanding}
\label{sec:procedural}

The third target task---\emph{sort the steps of a workflow into ascending
order}---is an instance of procedural understanding, a problem the NLP
community has studied under several guises: tracking state through a process,
inferring goal--step and step--step relations, inducing event schemas, and
planning. This section surveys those threads and connects them to ordering
the AEP/AJO tutorial steps from documentation and video transcripts.

\subsection{Procedural text understanding}

The canonical benchmark is \textbf{ProPara} \citep{proc:dalvi2018propara}:
$488$ human-authored science procedures with dense per-step annotations of
entity \emph{existence} and \emph{location}, where the model must track state
changes whose causes are often implicit and must be inferred---the same
implicit-causality challenge our pairwise primitive faces.
\textbf{OpenPI} \citep{proc:tandon2020proglobal} extends entity-state tracking
to open-domain wikiHow procedures. \textbf{RecipeQA}
\citep{proc:yagcioglu2018recipeqa} ($\sim$20K recipes, $\sim$36K QA pairs)
includes an explicit \emph{ordering} subtask---reorder shuffled steps---that is
structurally identical to our workflow-ordering goal, plus a coherence/cloze
subtask. The wikiHow line of work is most directly relevant:
\citet{proc:zhang2020wikihow} define \emph{Goal--Step Inference} (``learn
poses'' is a step of ``do yoga''), \emph{step inference}, and a \emph{temporal
ordering} task (``buy a yoga mat'' precedes ``learn poses''), reporting a
$10$--$20\%$ gap between transformer models and humans; \textbf{VGSI}
\citep{proc:yang2021vgsi} adds a visual variant over $772$K step images.
These step--step temporal relations are precisely the forward(1)/reverse(0)
precedence labels the repo derives from positional order within a document.

\subsection{Script and event-schema induction}

Ordering steps is a modern form of \emph{script} learning.
\citet{proc:chambers2008narrative} introduced unsupervised \emph{narrative
event chains} and the \emph{narrative cloze} evaluation---predict the held-out
event in a partially ordered chain---which directly anticipates next-step
prediction over a workflow. \textbf{ROCStories}
\citep{proc:mostafazadeh2016story} supplies a story-cloze/temporal-ordering
benchmark over everyday event sequences. Commonsense event schemas were
operationalized by \textbf{ATOMIC} \citep{proc:sap2019atomic}
(if--then social/physical relations incl.\ \emph{xNeed} prerequisites and
\emph{xEffect} consequences) and made generative by \textbf{COMET}
\citep{proc:bosselut2019comet}, which fine-tunes a transformer to emit
commonsense tails---a template for emitting prerequisite/consequence edges
between AEP steps. \textbf{proScript} \citep{proc:sakaguchi2021proscript} is the
most aligned: it generates \emph{partially ordered} scripts and defines an
\emph{edge-prediction} task (given a scenario and unordered events, recover the
partial-order DAG; F1$=75.7$). Its partial-order stance directly supports the
repo's diagnosis that forcing a \emph{total} order over parallel workflow steps
is a root cause of the $0/30$ ordering failure.

\subsection{Workflow and process modeling with LLMs}

Recent work scales procedural understanding to executable workflows.
\textbf{WorkflowLLM} \citep{proc:fan2024workflowllm} builds
\emph{WorkflowBench}---$106{,}763$ orchestration samples over $1{,}503$ APIs
from $83$ applications---and fine-tunes \emph{WorkflowLlama} to generate API
workflows, outperforming GPT-4o with in-context learning and generalizing to
out-of-distribution tasks (F1 plan $77.5\%$ on T-Eval). The repo's
\texttt{workflowllm\_aep\_merged\_2} dataset ($210$K pairs) adapts this corpus
to AEP. \textbf{FlowBench} \citep{proc:xiao2024flowbench} benchmarks
\emph{workflow-guided} agent planning across $6$ domains, $22$ roles, and $51$
scenarios, comparing workflow-knowledge formats (natural language, symbolic
code, flowchart) and finding even GPT-4o reaches only $\sim$40--43\% success---a
flowchart/DAG representation gave the best trade-off, reinforcing a
graph-structured rather than linear view of process knowledge. This connects to
classical \emph{process extraction} from documents and to hierarchical task
networks (HTNs) / task graphs, where a procedure is a partial order of subtasks
rather than a flat chain---the representation our rank-aggregation stage should
target.

\subsection{Procedural planning with LLMs}

The planning literature studies \emph{generating} an ordered plan rather than
sorting a given one, but the ordering competence transfers.
\citet{proc:huang2022zeroshotplanners} show large LMs decompose high-level
goals into mid-level step sequences zero-shot, though raw outputs need grounding
to admissible actions; \textbf{SayCan} \citep{proc:ahn2022saycan} grounds such
plans via affordance scoring. \textbf{Chain-of-thought} prompting
\citep{proc:wei2022cot} improves multi-step reasoning generally and underlies
the rationale-before-verdict strategy proposed for our reranker. Critically,
\textbf{CaT-Bench} \citep{proc:lyu2024catbench} benchmarks LM understanding of
\emph{causal and temporal} dependencies in plans (``can step $j$ happen before
step $i$?''), finding models conflate temporal sequence with causal necessity---
exactly the conflation the repo flags between \emph{narration order} and
\emph{causal precedence}.

\paragraph{Relevance to causal embedding for AEP/AJO.}
The downstream goal is to order the $30$ three-step AJO orchestrated workflows
(\texttt{ajo\_orchestrated\_workflows\_flat.json}, fields \texttt{t1/t2/t3}) and
to recover step order from the $142$ AEP video transcripts
(\texttt{aep\_transcripts.json}), where narration order is a noisy proxy for
causal order. The surveyed literature suggests four moves, consistent with the
repo's own root-cause analysis. (1) Adopt a \emph{partial-order/DAG} target
(proScript, FlowBench's flowchart format) instead of a forced total order, with
a NEUTRAL/abstain class for parallel steps such as the AJO welcome-flow ``pause
24h'' that is independent of audience enrichment. (2) Frame pairwise scoring as
prerequisite/consequence edge prediction (ATOMIC \emph{xNeed}/\emph{xEffect},
COMET) so the embedding distinguishes causal necessity from mere sequence, the
distinction CaT-Bench shows models miss. (3) Train on the \emph{transitive
closure} of within-document orders (not just adjacent pairs) to close the
distribution shift the repo identifies between gap$\le$2 training and
all-pairs evaluation. (4) Aggregate the (possibly cyclic) pairwise scores with a
robust ranker---minimum-feedback-arc-set / Bradley--Terry---feeding
\texttt{order\_events.py}, rather than a naive topological sort that lets per-pair
errors compound into the current $0/30$ result.
